% q0019.tex — Selic under fire (Monetary Policy / Central Banking)
\question{
  id        = {0019},
  title     = {Selic Under Fire},
  source    = {OBECON},
  year      = {2026},
  round     = {Seletiva IEO -- Phase A},
  language  = {English},
  topic     = {Macro},
  subtopic  = {Monetary Policy / Central Banking},
  type      = {objective},
  statement = {%
    During the third term of President Luiz Inácio Lula da Silva, the Central
    Bank of Brazil has been heavily criticized for maintaining the Selic rate
    at high levels even amid slowing inflation. Regarding the Central Bank,
    choose the correct alternative:
  },
  choices   = {%
    \item The President of the Central Bank of Brazil is Fernando Haddad,
          appointed by the President of the Republic and approved by the
          Senate to serve a four-year term, which does not coincide with the
          presidential term.
    \item Like the Federal Reserve, the Brazilian Central Bank's main
          objective is to keep inflation within a target set by the National
          Monetary Council and unemployment at the level of full employment.
    \item One of the recent advances in the management of the Brazilian
          Central Bank's activities was the establishment of its autonomy.
          According to economic theory, this measure helps increase the
          institution's credibility and, consequently, reduces the cost (in
          terms of output and unemployment) of fighting inflation.
    \item The Central Bank is responsible in Brazil for implementing fiscal
          policy, which includes setting taxation targets, managing public
          investments, and choosing the interest rate.
    \item New Keynesian economists often represent the Central Bank's reaction
          function to inflation and the output gap through the Phillips Curve.
  },
  answer    = {},
  solution  = {%
    % TODO: add solution
  },
}
